\section{Calibration and Data}
\label{sec:calibration}

\paragraph{The economy.}
The macro side is the \texttt{ogphl} multi-sector calibration of \OGCore{}: $\Mind=8$
industries and $\Igood=5$ consumption goods, built from the country package's own
production and consumption structure (its \texttt{PROD\_DICT}/\texttt{CONS\_DICT} and
social accounting matrix). The link does not hard-code this structure: it is
\emph{discovered} per run in the \OGCore{} environment and exported to the link through a
\texttt{baseline\_meta.json} concordance, which records how many industries and goods the
country resolves and which of them is electricity---here industry $m_e=2$ and energy good
$i_e=1$. Where a country's aggregation cannot isolate electricity as its own industry and
good, the energy channels report a skip rather than mis-target (\Cref{sec:signals-wedges}).
Production is Cobb--Douglas ($\varepsilon=1$), with the public-capital share
$\gamma_{g}$ carved out of each industry's capital share so that labour's residual share
is well defined.

\paragraph{The electricity-cost proxy.}
The ``auto'' source is a curated cost-of-electricity index where available, otherwise
the levelized cost of electricity (LCOE) reconstructed from the run's
per-technology costs and generation (the commodity-balance shadow price is degenerate,
binding in only scattered years, so it is not used)---taken as a reform/base ratio so the
(uncalibrated) money units cancel. This ratio is an average-cost proxy, not an observed
retail or marginal price. For the Philippine reform it averages about $1.05$ and reaches
$1.16$ around 2030. The same proportional change is applied to disjoint expenditure
bases through the composite of \Cref{sec:routes}: an inter-industry cost-push weighted by
electricity input-cost shares from the SAM, and a recycled household wedge on final
electricity use. The electricity industry's self-purchases are removed before upstream
cost propagation. The direct and indirect cost-push follows the input--output logic of
\Cref{eq:routeB}; the
specific pass-through magnitudes are calibration-dependent and re-derived with the SAM
(the earlier M=4 figures are retired with that aggregation).

\paragraph{Health.}
The health channel uses sourced epidemiological estimates: the ambient-PM\textsubscript{2.5}
attributable burden and its age profile from the IHME Global Burden of Disease export
(covering several countries; for the Philippines $43{,}951$ attributable deaths, with
deaths-by-age and years-lived-with-disability-by-age profiles), scaled by a calibrated
dose-response $M\approx0.082$ (an energy mass-share $\times$ concentration-response
elasticity). The lives effect is
$\Delta\text{deaths}=\text{total}\times M\times\Delta\text{emissions}$, so the reform's
$-2.7\%$ emissions change (the first-decade mean) saves about $95$ lives (the $\times M$
dose-response separates
this from the naive $43{,}951\times\Delta\text{emissions}\approx-1{,}165$). The GBD burden
and age profile are sourced inputs. The construction of $M$, the linearity of the
emissions-to-burden mapping, and the morbidity productivity elasticity remain calibrated
model choices.

\paragraph{Units.}
The one uncalibrated bridge is the deflator that maps \CLEWS{} money to the \OGCore{} real
numéraire, currently set to unity. It affects the \emph{levels} of the carbon and
investment channels (and any revenue figure), though not the ratio-based energy price or
the demographic health inputs.

\caveat{Magnitudes for the carbon and investment channels are illustrative until the
unit/deflator bridge is calibrated; the transformation and solve mechanics are tested
(\Cref{sec:results}) and the \emph{sizes} scale with the bridge. The electricity input is
ratio-based, while the health burden and age profile are sourced estimates with a
calibrated emissions-to-burden translation.}
